\documentclass[10pt,a4paper]{amsart}
\usepackage[margin=19mm]{geometry}
\usepackage{amsmath,amssymb,mathtools}
\usepackage[scr=rsfs,cal=euler]{mathalfa}
\usepackage{xfrac}
\usepackage[unicode,hidelinks,bookmarks=false]{hyperref}
\newcommand{\ie}{i.e.\ }
\newcommand{\cf}{cf.\ }
\newcommand{\resp}{respectively\ }
\newcommand{\Subsection}{\subsection}
\providecommand{\nobreakdash}{\nobreak\hbox{-}}
\setlength{\emergencystretch}{2em}
\multlinegap0pt
\usepackage{tikz}
\usepackage{amssymb}
\usepackage{bm}
\usepackage[shortlabels]{enumitem}
\usepackage{mathtools}
\newtheorem{cor}{Corollary}
\def\cH{{\mathcal H}}
\title{Covering half-grids with lines and planes}
\author{Anurag Bishnoi, Shantanu Nene}
\date{Source: arXiv:2501.11156v3 (CC BY 4.0)}
\begin{document}
\maketitle

\noindent\emph{Source and licence.} Covering half-grids with lines and planes. Version arXiv:2501.11156v3 (CC BY 4.0), distributed by arXiv under the Creative Commons Attribution 4.0 International licence, http://creativecommons.org/licenses/by/4.0/, as recorded in the arXiv licence metadata for this exact version and shown on the abstract page https://arxiv.org/abs/2501.11156v3. Full licence notice: \texttt{licenses/arxiv-2501.11156-CC-BY-4.0.txt}.\par\smallskip

\noindent\textit{Editorial scope.} Counted unit: the article's cor cor (source0.tex lines 244--246). The article's complete original proof of it is delivered with it (source0.tex lines 247--417, one \texttt{proof} environment of 171 lines). No mathematics is written, repaired or re-derived: the statement and the proof are the article's own lines, in the article's own notation, and no retained line is substituted or re-flowed.  No further source-local context is needed: every cross-reference of every retained line resolves inside the two delivered spans.  Nothing was omitted from any retained span, so the package carries no omission record.  No retained line contains a citation, so the wrapper carries no bibliography and no bibliography entry.  No retained line contains a figure environment, an image-inclusion command or drawing code, so no image is delivered and the package declares no asset dependency.  Editorial formatting only, in addition: an AMS \texttt{amsart} preamble that restates the article's own declarations and macros used by the retained lines, namely the environments \texttt{cor} on separate counters, together with the standard packages those lines need.  The retained environments print in this document as Corollary 1 in order of appearance, whereas the article prints the same items with the numbering of its own full document.  The complete native source of the article as distributed in the arXiv e-print for this exact version is bundled unchanged as \texttt{sources/171812/source0.tex}.
\begin{cor}
    Let $\cH$ be an $m \times n$ half-grid, with $m \leq n$. Then the minimum number of lines required to cover every point of $\cH$ at least $k$ times is at least $mk\left(1-e^{-\frac{n}{m}}-O(\frac{n}{m^2})\right)$.
\end{cor}

\begin{proof}
    Note that in the case of $n=m$, $\cH$ is a conical grid, and the bound reduces to the same one as the previous case. Hence from now on, we assume $n \geq m+1$. 
    
    Let the sets of $\cH$ be $S_1=\{c_0<c_1<\cdots<c_{n-1}\}$ and $S_2=\{d_0<d_1<\cdots<d_{m-1}\}$. Then, using the notation of the lemma, $\alpha_i=d_{m-i}$ for $1 \leq i \leq m$, and 
    $$\beta_i=1+\left\lfloor\frac{(i-1)(n-1)}{m-1}\right\rfloor > \frac{(i-1)(n-1)}{m-1}.$$ 
    Let $l=l_0$ for convenience. If $l \geq mk$, then we are already done, so we may assume $l < mk$.
 
    Define $s = \lceil \frac{l(m-1)}{k(n-1)} \rceil$ and $t = 1+s$. 
    Since $l < mk$, we have $s \leq \lceil \frac{m(m-1)}{n-1} \rceil \leq m-1$, so $t \leq m$ and $\beta_t$ is well-defined.
    Moreover, $\frac{(t-1)(n-1)}{m-1} = \frac{s(n-1)}{m-1} \geq \frac{l}{k}$, so the condition $\beta_t \geq l/k$ of the lemma is satisfied.
 
    By the lemma, the number of lines used is at least
    $$l+(m-s)k - l\left(\frac{1}{\beta_{s+1}}+\cdots + \frac{1}{\beta_m}\right).$$
    We first handle the case $s=0$. This can only happen if $l=0$, and in that case the above bound becomes $mk$, so we are done. \\
    
    Since $\beta_i > \frac{(i-1)(n-1)}{m-1}$, we have $\frac{1}{\beta_i} < \frac{m-1}{(i-1)(n-1)}$, giving
    \begin{align}
    l+(m-s)k-l\left(\frac{1}{\beta_{s+1}}+\cdots + \frac{1}{\beta_m}\right) & > l+(m-s)k-\frac{l(m-1)}{n-1} \left( \frac{1}{s}+\frac{1}{s+1}+\cdots+\frac{1}{m-1}\right) \notag\\
    &= (m-s)k+l\left(1-\frac{(m-1)(H_{m-1}-H_{s-1})}{n-1}\right). \label{eq:main_expr}
    \end{align}

    Next suppose $s=1$. Then \eqref{eq:main_expr} gives the lower bound
    \[
    (m-1)k+
    l\left(1-\frac{(m-1)H_{m-1}}{n-1}\right).
    \]
    If the coefficient of $l$ is non-negative, this is at least
    $(m-1)k=mk(1-1/m)$, which is enough since $n\geq m$ and hence
    $1/m=O(n/m^2)$.

    It remains to consider the case where the coefficient of $l$ is negative.
    Then $n-1<(m-1)H_{m-1}$. Since $s=1$, we have
    $l\leq k(n-1)/(m-1)$, and therefore
    \begin{align*}
    (m-1)k+
    l\left(1-\frac{(m-1)H_{m-1}}{n-1}\right)
    &\geq
    (m-1)k+
    \frac{k(n-1)}{m-1}
    \left(1-\frac{(m-1)H_{m-1}}{n-1}\right)  \\
    &=
    (m-1)k+\frac{k(n-1)}{m-1}-kH_{m-1}
    \geq
    mk-kH_{m-1}.
    \end{align*}
    We now use the elementary estimate
    $H_{m-1}/m\leq e^{-n/m}+O(n/m^2)$. Indeed, if
    $n/m\leq \frac12\log m$, then $e^{-n/m}\geq m^{-1/2}$ dominates
    $H_{m-1}/m=O(\log m/m)$; if $n/m>\frac12\log m$, then
    $n/m^2\gg \log m/m$. Hence
    \[
    mk-kH_{m-1}
    =
    mk\left(1-\frac{H_{m-1}}{m}\right)
    \geq
    mk\left(1-e^{-n/m}-O\left(\frac{n}{m^2}\right)\right).
    \]
    
        For the remainder of the proof, assume $s\geq 2$. Put
    $A=\frac{n-1}{m-1}$ and
    \[
    D(s)=H_{m-1}-\ln(s-1)-\gamma.
    \]
    Since $H_{s-1}>\ln(s-1)+\gamma$, \eqref{eq:main_expr} is greater than
    \begin{equation}\label{eq:f_bound}
    (m-s)k+l\left(1-\frac{D(s)}{A}\right).
    \end{equation}
    We repeatedly use
    \[
    \frac{m-1}{m}e^{\frac{1}{2(m-1)}-\frac{n-1}{m-1}}
    =
    e^{-n/m+O(n/m^2)},
    \]
    and the fact that $O(k)$ is absorbed into $mk\cdot O(n/m^2)$ since
    $n\geq m$.

    \textbf{Case I:} $D(s)\leq A$. The coefficient of $l$ in
    \eqref{eq:f_bound} is non-negative, so using
    $l>\frac{(s-1)k(n-1)}{m-1}=(s-1)kA$, the expression is greater than
    \[
    f(s):=(m-s)k+(s-1)k(A-D(s)).
    \]
    Since $f'(s)=k(A-D(s))\geq 0$ throughout this case, $f$ is minimized at
    the left boundary $s_0=1+e^{H_{m-1}-\gamma-A}$, where $D(s_0)=A$. Hence
    \[
    f(s)\geq f(s_0)=mk-s_0k.
    \]
    Using $H_{m-1}<\ln(m-1)+\gamma+\frac{1}{2(m-1)}$, we get
    \[
    s_0<1+(m-1)e^{\frac{1}{2(m-1)}-\frac{n-1}{m-1}}.
    \]
    Therefore
    \[
    f(s)
    \geq
    mk-k-k(m-1)e^{\frac{1}{2(m-1)}-\frac{n-1}{m-1}}
    \geq
    mk\left(1-e^{-n/m}-O\left(\frac{n}{m^2}\right)\right).
    \]

    \textbf{Case II:} $D(s)>A$. The coefficient of $l$ in
    \eqref{eq:f_bound} is negative, so using
    $l\leq \frac{sk(n-1)}{m-1}=skA$, the expression is at least
    \[
    h(s):=(m-s)k+sk(A-D(s))
    =
    mk+sk(A-1-D(s)).
    \]
    We have $h''(s)=\frac{k(s-2)}{(s-1)^2}\geq 0$, so it suffices to check
    the possible boundary points and the possible interior critical point.

    If $s^*$ is an interior critical point, then
    $D(s^*)=A+\frac{1}{s^*-1}$, and hence
    \[
    h(s^*)=
    mk-\frac{(s^*)^2k}{s^*-1}
    =
    mk-k(s^*-1)-2k-\frac{k}{s^*-1}.
    \]
    Since
    \[
    s^*-1<e^{H_{m-1}-\gamma-A}
    <
    (m-1)e^{\frac{1}{2(m-1)}-\frac{n-1}{m-1}},
    \]
    we get
    \[
    h(s^*)\geq
    mk-k(m-1)e^{\frac{1}{2(m-1)}-\frac{n-1}{m-1}}-O(k)
    \geq
    mk\left(1-e^{-n/m}-O\left(\frac{n}{m^2}\right)\right).
    \]

    At the boundary $s=s_0$, if it occurs, we have $D(s_0)=A$, so
    $h(s_0)=mk-s_0k$. The same estimate as in Case~I gives
    \[
    h(s_0)\geq
    mk\left(1-e^{-n/m}-O\left(\frac{n}{m^2}\right)\right).
    \]

    At the endpoint $s=2$, if it occurs, then
    $A\leq H_{m-1}-\gamma=O(\log m)$, so $n/m=O(\log m)$. Also
    \[
    h(2)=
    mk+
    2k\left(\frac{n-m}{m-1}-(H_{m-1}-\gamma)\right)
    \geq
    mk-O(k\log m).
    \]
    Since, in this regime,
    $\log m/m\leq e^{-n/m}+O(n/m^2)$, this is enough.

    Finally, at the endpoint $s=m-1$, if it occurs, then
    \[
    h(m-1)
    =
    nk-(m-1)k\left(H_{m-1}-\ln(m-2)-\gamma\right)
    =
    nk-O(k).
    \]
    This is at least $mk-O(k)$, and hence again at least
    \[
    mk\left(1-e^{-n/m}-O\left(\frac{n}{m^2}\right)\right).
    \]

    Thus in every case the number of lines is at least
    \[
    mk\left(1-e^{-n/m}-O\left(\frac{n}{m^2}\right)\right),
    \]
    as required.
\end{proof}

\end{document}
